\centerline{\zihao{-3}\heiti 中文摘要}

\vspace{1em}

普惠金融的深入推进推动小微贷款行业快速发展，但服务质量参差不齐已成为制约行业健康发展的突出瓶颈。P公司作为经地方金融监管部门批准的普惠金融公司，同样面临客户满意度持续低迷、投诉处理效率偏低、业务流程一次通过率不高等服务质量问题。如何运用系统化方法精准诊断服务质量短板并实施有效改进，是P公司亟需解决的现实命题。

本研究以DMAIC方法论为主线，构建了适用于小微贷款场景的服务质量改进框架。在界定（Define）阶段，识别关键质量特性并确立改进目标；在测量（Measure）阶段，基于SERVQUAL经典五维度扩展为包含便利性、透明性与产品适配的八维度评价模型，设计Likert 5级问卷，回收有效问卷486份（Cronbach's $\alpha$ = 0.916，KMO = 0.894），信效度良好；在分析（Analyze）阶段，综合运用SERVQUAL差距分析与重要性—绩效分析（IPA），发现透明性维度感知—期望差距最大（$-$1.70），IPA优先改进区包含6项关键指标；进一步通过层次分析法（AHP）对5位行业专家进行结构化判断（CR = 0.038，一致性良好），量化确定改进优先级；在改进（Improve）阶段，据此制定并实施5项针对性改进方案；在控制（Control）阶段，建立关键指标监控机制以固化改进成果。

试点结果表明：客户综合满意度由3.12提升至3.85，投诉闭环率由62\%提升至88\%，业务一次通过率由62\%提升至78\%，平均处理时长由12个工作日缩短至7.5个工作日，各项服务质量指标均有明显提升。研究表明，DMAIC方法适用于小微贷款服务质量改进，所构建的八维评价模型对小微信贷场景具有较好的适用性，为同类机构提供了可借鉴的改进路径。

\vspace{1em}
\noindent{\zihao{4}\heiti 关键词：}服务质量；DMAIC；小微贷款；SERVQUAL；普惠金融；IPA；层次分析法
